(a) A germ $s_P\in\mathscr F_P$ belongs to $\ker\varphi_P$ exactly when a representative $s\in\mathscr F(U)$ satisfies $\varphi(s)|_V=0$ on some neighborhood $V$ of $P$. Then $s|_V\in(\ker\varphi)(V)$. This gives $(\ker\varphi)_P=\ker\varphi_P$. Sheafification preserves stalks, and the germs of sections of the image presheaf are exactly the germs $\varphi(s)_P=\varphi_P(s_P)$. Thus $(\operatorname{im}\varphi)_P=\operatorname{im}\varphi_P$.

(b) A sheaf is zero exactly when all its stalks vanish. Hence $\ker\varphi=0$ exactly when every $\varphi_P$ is injective. Likewise $\operatorname{im}\varphi\to\mathscr G$ is an isomorphism exactly when it is an isomorphism on every stalk, which is equivalent to every $\varphi_P$ being surjective.

(c) A composite of sheaf morphisms is zero if and only if all its stalk maps are zero. Assuming this, the inclusion of the image of one map into the kernel of the next is an isomorphism if and only if it is an isomorphism on all stalks. Apply (a) at each term of the sequence.
